\seccion{2.j. Cliente con el mayor número de órdenes}

Se obtiene la información del cliente que haya realizado el mayor número de órdenes, junto con la cantidad de órdenes que efectuó. Primero se proyectan los identificadores de cliente y orden para eliminar posibles repeticiones. Después, se agrupan las órdenes por cliente y se cuenta el número de órdenes realizadas por cada uno. Finalmente, se calcula el máximo, se seleccionan los clientes que alcanzan ese valor y se relaciona el resultado con \texttt{customer} para recuperar toda su información.

Si existe un empate, se muestran todos los clientes que hayan realizado el máximo número de órdenes.

\textbf{Expresión en álgebra relacional}

\begin{align*}
O &\leftarrow \pi_{customerid,orderid}(orders)\\
N &\leftarrow \gamma_{customerid;\ \operatorname{count}(orderid)\to numorders}(O)\\
M &\leftarrow \gamma_{\operatorname{max}(numorders)\to maxorders}(N)\\
T &\leftarrow \sigma_{numorders=maxorders}(N\times M)\\
J &\leftarrow \pi_{customer.*,\,numorders}(customer\bowtie T).
\end{align*}

\textbf{Expresión ejecutada en RelaX}

\begin{lstlisting}
O = π customerid, orderid (orders)

N = γ customerid; count(orderid) → numorders (O)

M = γ max(numorders) → maxorders (N)

T = σ numorders = maxorders (N ⨯ M)

J = π customer.*, numorders (customer ⨝ T)

J
\end{lstlisting}

\textbf{Resultado: 1 cliente, con 13 órdenes.} No hubo empate; por lo tanto, se obtuvo un único cliente con la cantidad máxima de órdenes.

\begin{center}
\small
\begin{tabular}{llll}
\toprule
\rowcolor{Fondo}
Atributo & Valor & Atributo & Valor\\
\midrule
CustomerID & JE-15745 & Nombre & Joel Eaton\\
Segmento & Consumer & País & United States\\
Ciudad & Jackson & Estado & Michigan\\
Código postal & 49201 & Región & Central\\
\multicolumn{3}{r}{\textbf{Número de órdenes}} & \textbf{13}\\
\bottomrule
\end{tabular}
\end{center}

El resultado completo se encuentra en \texttt{datos/consulta\_j\_resultado.csv}.

\begin{center}
\begin{minipage}{.48\linewidth}
\caprel{consulta_j_arbol.png}{8.0}{Árbol lógico de la consulta 2.j en RelaX.}
\end{minipage}
\hfill
\begin{minipage}{.48\linewidth}
\caprel{consulta_j_resultado.png}{8.0}{Resultado de la consulta 2.j en RelaX.}
\end{minipage}
\end{center}